\chapter{Marco teórico}
\label{ch:marco}
\section{Tareas, archivos y jobs}
Un flujo de trabajo científico puede representarse como un grafo dirigido acíclico. Cada tarea produce archivos que consumen las tareas sucesoras. Las entradas provienen de fuera del grafo; los resultados finales deben quedar disponibles en el destino declarado. Los archivos que conectan tareas son intermedios.

Pegasus transforma el flujo científico en un plan ejecutable con jobs científicos y operaciones auxiliares \citep{deelman2015}. Cada job contiene una tarea o un grupo. DAGMan (\textit{Directed Acyclic Graph Manager}) respeta las dependencias entre jobs \citep{htcondordagman}. HTCondor asigna cada job a recursos elegibles y registra su ejecución \citep{htcondorintro}. El plan muestra qué tareas se agruparon; el historial, dónde se ejecutaron.

\section{Ubicación y entrega de archivos}
En CondorIO, el coordinador prepara las entradas y recibe las salidas. Cada job trabaja en un directorio aislado o \textit{sandbox}. Un archivo puede volver al coordinador y salir hacia otro job, incluso si ambos corren en la misma máquina \citep{htcondorordinary,pegasustransfers}. Al agrupar tareas consecutivas, algunas lecturas pueden resolverse en el mismo sandbox. Como las salidas declaradas pueden conservarse, los registros permiten comprobar el ahorro.

El volumen observado suma el tamaño de cada entrega completa; un mismo archivo puede generar varias entregas. Los contadores nativos también pueden incluir ejecutables, metadatos y registros auxiliares. Por eso se conservan los archivos y tamaños científicos, el contrato efectivo de cada job, sus estadísticas y eventos. Kickstart y los metadatos de Pegasus aportan información de procedencia \citep{pegasusprovenance}.

\section{Patrones de flujo}
Los patrones representan relaciones de producción y consumo, y permiten aislar cambios en el movimiento de archivos. La caracterización de workflows científicos identifica estructuras y dependencias diversas \citep{bharathi2008}. Aquí, procesamiento sirve como control; pipeline sigue dependencias consecutivas; agregación reúne salidas; distribución divide un conjunto; y redistribución combina reunión y reparto.
En las Figuras~\ref{fig:pat-process} a~\ref{fig:pat-redistribution}, los círculos son tareas y los rectángulos archivos; las flechas indican producción y consumo, no transferencias entre máquinas. La ubicación y los contratos se incorporan después al cálculo. Las condiciones comparadas se describen en la sección~\ref{sec:repro_patrones}.

\subsection{Procesamiento}
Una tarea transforma una entrada en un resultado final. La Figura~\ref{fig:pat-process} presenta el caso utilizado en P01 y P02.

\par\medskip\noindent\begin{minipage}{\linewidth}
\captionsetup{type=figure}
\caption{Patrón de procesamiento}
\label{fig:pat-process}
\centering
\begin{lrbox}{\dmediagrambox}\begin{tikzpicture}[>=Stealth,font=\small,
 t/.style={draw,circle,minimum size=0.75cm,inner sep=1pt},
 f/.style={draw,rectangle,minimum width=1cm,minimum height=0.6cm,inner sep=2pt}]
\node[f] (a) at (0,0) {$A$};
\node[t] (t) at (3,0) {$T_1$};
\node[f] (b) at (6,0) {$B$};
\draw[->] (a)--(t); \draw[->] (t)--(b);
\end{tikzpicture}\end{lrbox}\fitdmediagram
\par\smallskip\raggedright\normalsize\textit{Nota.} $A$ es la entrada externa y $B$ el resultado final. Ambos miden 10 MiB en P01 y 20 MiB en P02. La estructura permanece constante al cambiar el tamaño.
\end{minipage}\par\medskip

\subsection{Pipeline}
Cada etapa consume la salida de la anterior. La Figura~\ref{fig:pat-pipeline} conserva la identidad de los intermediarios, aunque las tareas se agrupen después en jobs.

\par\medskip\noindent\begin{minipage}{\linewidth}
\captionsetup{type=figure}
\caption{Patrón pipeline y sus archivos intermedios}
\label{fig:pat-pipeline}
\centering
\begin{lrbox}{\dmediagrambox}\begin{tikzpicture}[>=Stealth,font=\small,
 t/.style={draw,circle,minimum size=0.75cm,inner sep=1pt},
 f/.style={draw,rectangle,minimum width=1cm,minimum height=0.6cm,inner sep=2pt}]
\node[f] (a) at (0,0) {$A$};
\node[t] (t1) at (2,0) {$T_1$};
\node[f] (f12) at (4,0) {$f_{12}$};
\node[t] (t2) at (6,0) {$T_2$};
\node[f] (f23) at (8,0) {$f_{23}$};
\node[t] (t3) at (10,0) {$T_3$};
\node[f] (b) at (12,0) {$B$};
\foreach \u/\v in {a/t1,t1/f12,f12/t2,t2/f23,f23/t3,t3/b}{\draw[->] (\u)--(\v);}
\end{tikzpicture}\end{lrbox}\fitdmediagram
\par\smallskip\raggedright\normalsize\textit{Nota.} En L01, $A$, $f_{12}$, $f_{23}$ y $B$ miden 10 MiB. Agrupar tareas conserva estos archivos y sus dependencias, aunque puede modificar las entregas del mecanismo de ejecución.
\end{minipage}\par\medskip

\Needspace{16\baselineskip}
\subsection{Agregación}
Una tarea agregadora consume archivos distintos de varios productores. La Figura~\ref{fig:pat-aggregation} muestra los cuatro productores del caso base.

\par\medskip\noindent\begin{minipage}{\linewidth}
\captionsetup{type=figure}
\caption{Patrón de agregación: varios productores alimentan una tarea agregadora}
\label{fig:pat-aggregation}
\centering
\begin{lrbox}{\dmediagrambox}\begin{tikzpicture}[>=Stealth,font=\small,
 t/.style={draw,circle,minimum size=0.75cm,inner sep=1pt},
 f/.style={draw,rectangle,minimum width=1cm,minimum height=0.6cm,inner sep=2pt}]
\node[t] (agg) at (7,0) {$G$};
\node[f] (b) at (10,0) {$B$};
\foreach \i/\y in {1/1.8,2/0.6,3/-0.6,4/-1.8}{
 \node[f] (a\i) at (0,\y) {$A_{\i}$};
 \node[t] (p\i) at (2,\y) {$P_{\i}$};
 \node[f] (f\i) at (4,\y) {$f_{\i}$};
 \draw[->] (a\i)--(p\i); \draw[->] (p\i)--(f\i); \draw[->] (f\i)--(agg);
}
\draw[->] (agg)--(b);
\end{tikzpicture}\end{lrbox}\fitdmediagram
\par\smallskip\raggedright\normalsize\textit{Nota.} En A01, cada entrada $A_i$ y cada intermediario $f_i$ mide 10 MiB; $G$ produce el final $B$ de 40 MiB. Cada subíndice identifica un archivo distinto.
\end{minipage}\par\medskip

\subsection{Distribución}
Un productor puede repartir particiones distintas o compartir un archivo entre varios consumidores. La Figura~\ref{fig:pat-distribution} distingue ambas relaciones, necesarias para interpretar el efecto de agrupar.

\par\medskip\noindent\begin{minipage}{\linewidth}
\captionsetup{type=figure}
\caption{Patrón de distribución: particiones distintas y archivo compartido}
\label{fig:pat-distribution}
\centering
\begin{lrbox}{\dmediagrambox}\begin{tikzpicture}[>=Stealth,font=\small,
 t/.style={draw,circle,minimum size=0.7cm,inner sep=1pt},
 f/.style={draw,rectangle,minimum width=0.95cm,minimum height=0.55cm,inner sep=2pt}]
\node[anchor=west] at (-0.5,2.5) {(a) Particiones distintas};
\node[f] (a) at (0,0) {$A$}; \node[t] (s) at (2,0) {$S$};
\draw[->] (a)--(s);
\foreach \i/\y in {1/1.65,2/0.55,3/-0.55,4/-1.65}{
 \node[f] (q\i) at (5,\y) {$q_{\i}$};
 \node[t] (c\i) at (8,\y) {$C_{\i}$};
 \node[f] (b\i) at (11,\y) {$B_{\i}$};
 \draw[->] (s)--(q\i); \draw[->] (q\i)--(c\i); \draw[->] (c\i)--(b\i);
}
\node[anchor=west] at (-0.5,-3) {(b) Archivo compartido};
\node[f] (aa) at (0,-5.5) {$A$}; \node[t] (p) at (2,-5.5) {$P$};
\node[f] (g) at (5,-5.5) {$g$};
\draw[->] (aa)--(p); \draw[->] (p)--(g);
\foreach \i/\y in {1/-3.85,2/-4.95,3/-6.05,4/-7.15}{
 \node[t] (cc\i) at (8,\y) {$C_{\i}$};
 \node[f] (bb\i) at (11,\y) {$B_{\i}$};
 \draw[->] (g)--(cc\i); \draw[->] (cc\i)--(bb\i);
}
\end{tikzpicture}\end{lrbox}\fitdmediagram
\par\smallskip\raggedright\normalsize\textit{Nota.} En D02 y GD4, $S$ divide una entrada de 10 MiB en cuatro particiones $q_i$ de 2,5 MiB y cada final $B_i$ conserva ese tamaño. En AC1, AC2 y AC4, $P$ produce un único archivo $g$ de 10 MiB, consumido por cuatro tareas que generan finales de 10 MiB. Las cuatro aristas de consumo de $g$ conservan una sola identidad de archivo.
\end{minipage}\par\medskip

\subsection{Redistribución}
La redistribución agrega archivos de varios productores y divide el resultado para nuevos consumidores. La Figura~\ref{fig:pat-redistribution} representa ambas operaciones y sus intermediarios.

\par\medskip\noindent\begin{minipage}{\linewidth}
\captionsetup{type=figure}
\caption{Patrón de redistribución: agregación seguida de partición y consumo}
\label{fig:pat-redistribution}
\centering
\begin{lrbox}{\dmediagrambox}\begin{tikzpicture}[>=Stealth,font=\small,
 t/.style={draw,circle,minimum size=0.65cm,inner sep=1pt},
 f/.style={draw,rectangle,minimum width=0.85cm,minimum height=0.55cm,inner sep=2pt}]
\node[t] (agg) at (5.25,0) {$G$};
\node[f] (h) at (7,0) {$h$};
\node[t] (s) at (8.75,0) {$S$};
\draw[->] (agg)--(h); \draw[->] (h)--(s);
\foreach \i/\y in {1/1.8,2/0.6,3/-0.6,4/-1.8}{
 \node[f] (a\i) at (0,\y) {$A_{\i}$};
 \node[t] (p\i) at (1.75,\y) {$P_{\i}$};
 \node[f] (f\i) at (3.5,\y) {$f_{\i}$};
 \draw[->] (a\i)--(p\i); \draw[->] (p\i)--(f\i); \draw[->] (f\i)--(agg);
 \node[f] (q\i) at (10.5,\y) {$q_{\i}$};
 \node[t] (c\i) at (12.25,\y) {$C_{\i}$};
 \node[f] (b\i) at (14,\y) {$B_{\i}$};
 \draw[->] (s)--(q\i); \draw[->] (q\i)--(c\i); \draw[->] (c\i)--(b\i);
}
\end{tikzpicture}\end{lrbox}\fitdmediagram
\par\smallskip\raggedright\normalsize\textit{Nota.} En R01-E, las entradas $A_i$, los intermediarios $f_i$, las particiones $q_i$ y los finales $B_i$ miden 10 MiB cada uno. El archivo agregado $h$ mide 40 MiB. Las tareas $G$ y $S$ representan agregación y división, respectivamente.
\end{minipage}\par\medskip


\section{Antecedentes y relación con DME}
\citet{chen2015} estudian la agrupación de tareas y sus efectos en la ejecución, una decisión que modifica el plan. Esta tesis pregunta qué entregas científicas cambian entre configuraciones y cómo se relacionan con las obligaciones del mismo grafo. Para responderla, se requiere el mapa de tareas a jobs y las listas de archivos transferidos.

Pietri y Sakellariou (\citeyear{pietri2018}) estudian la planificación para reducir comunicación. \citet{fortmeier2013} relacionan modelos de volumen con métricas de partición de hipergrafos. Aquí, $n_f$ cuenta, para cada archivo, los destinos distintos del origen y se extiende con el relevo CondorIO. No se adopta su métrica de comunicación entre todos los vecinos: se fija la ubicación observada y no se optimiza la partición.

La métrica MUE (Memory Usage Efficiency) usa una referencia teórica de accesos a memoria y el aprovechamiento del ancho de banda para evaluar una implementación \citep{fuhrer2018}. Este antecedente ilustra cómo una referencia ayuda a interpretar un volumen. DME opera a otra escala: define necesidades por archivo y ubicación y las contrasta con las entregas corroboradas durante una ejecución del flujo.

\texttt{pegasus-statistics} informa conteos y tiempos de tareas y jobs \citep{pegasusstats}, pero no expresa la cota por archivo, ubicación y relevo definida aquí. Analyzer añade esa referencia y concilia las entregas científicas. La integración de registros de distintas capas tiene antecedentes en \citet{papadimitriou2021}; para compararlos, las fuentes deben corresponder a una misma ejecución.

\begin{table}[htbp]\centering\small
\caption{Variables y fuentes de la caracterización}
\begin{tabularx}{\linewidth}{@{}P{2.75cm}YY@{}}
\toprule Variable & Fuente de referencia & Función en la medición\\\midrule
Rol del archivo & Grafo científico y declaración de finales. & Distinguir entrada, intermedio y final, y determinar productores y consumidores.\\
Identidad y tamaño & Catálogos y metadatos de ejecución. & Vincular cada archivo y expresarlo en bytes.\\
Tarea y job & Plan de Pegasus y descripción del grupo. & Identificar si dos tareas usan el mismo espacio de ejecución.\\
Máquina e intento & Historial y eventos de HTCondor. & Resolver ubicación efectiva, terminación y alcance del intento.\\
Contrato de entrega & Listas efectivas de entrada y salida. & Identificar qué archivos pueden contabilizarse en cada dirección.\\
Contadores y tiempos & Estadísticas de transferencia y métricas de DAGMan. & Conciliar bytes y conteos. Medir duración del workflow.\\
\bottomrule
\end{tabularx}
\end{table}

Distinguir el archivo lógico de cada copia entregada evita contar solo por nombre de tarea o job. En el pipeline, un intermedio puede ser local según la referencia si productor y consumidor están en W1. Si sus jobs usan sandboxes distintos, el mecanismo puede devolverlo al coordinador y entregarlo al siguiente job. El plan y los registros permiten establecer ese recorrido; el grafo y la ubicación, determinar si el dato hacía falta en otra máquina.
